
\documentclass{studrep}

\RequirePackage{tabularx}
% \RequirePackage{graphicx}

\graphicspath{{pics/}}

% \usepackage{showframe}

% pygmetize
\RequirePackage{minted}
\usemintedstyle{tango} % bw

% \setmainfont{Times New Roman}

\begin{document}
\thispagestyle{empty}
\begin{center}
Министерство науки и высшего образования Российской Федерации

Федеральное государственное бюджетное образовательное учреждение\\ высшего образования

«Иркутский государственный университет»\\
(ФГБОУ ВО «ИГУ»)

Институт математики и информационных технологий

Кафедра информационных технологий
\end{center}

\vfill
\begin{center}
  \textbf{\large ОТЧЕТ}

  по научно"=исследовательской работе
  % по производственной практике
  % по преддипломной практике   % и т.п.
  %\textbf{ ВЫПУСКНАЯ КВАЛИФИКАЦИОННАЯ РАБОТА\\
%БАКАЛАВРА}
\vspace{1em}

% по направлению 02.03.03 -- Математическое обеспечение и администрирование информационных систем

% профиль <<общий>>

\vspace{2em}
ТЕМА БОЛЬШИМИ БУКВАМИ

\end{center}
\vfill

\noindent\begin{tabularx}{\textwidth} {
  >{\raggedright\arraybackslash}X
  >{\raggedright}X }
&

Студента \_ курса очного отделения\\
группы 02\_4\_-ДБ\\
Фамилия Имя Отчество\\[2em]

Руководитель:\\
к.~т.~н., доцент\\
\underline{\hspace{3cm}} Иванов Иван Иванович\\[2em]

Защищен с оценкой\\[1em] \underline{\hspace{3cm}}
%Допущена к защите\\
%Зав.каф., к.~т.~н., доцент\\
%\underline{\hspace{3cm}} Черкашин Евгений\\ \hspace{3cm} Александрович\\[2em]

\end{tabularx}
\vfill
\begin{center}
  Иркутск -- 2022
\end{center}
\clearpage

\tableofcontents

\chapter*{ВВЕДЕНИЕ}
% TODO: Add contents line
\label{chap:intro}

\chapter{Теоретические основы \ldots}
% Другой пример названия
%\chapter{Математическое обеспечение и технологии разработки ... \ldots}

Согласно \cite{bratko90}

.... база данных (БД) .....


\chapter{Исследование скомпилированного кода}

\begin{minted}{text}
  stud@sysrescue:~/projects/webapp$ dub run --compiler=ldc2
    Starting Performing "debug" build using ldc2 for x86_64.
  Up-to-date hunt 1.7.17: target for configuration [library] is up to date.
  Up-to-date hunt-extra 1.2.3: target for configuration [library] is up to date.
  Up-to-date hunt-net 0.7.1: target for configuration [default] is up to date.
  Up-to-date hunt-http 0.8.2: target for configuration [default] is up to date.
  Up-to-date protobuf 0.6.2: target for configuration [protobuf] is up to date.
  Up-to-date grpc 0.5.0-beta.2: target for configuration [library] is up to date.
  Up-to-date hunt-redis 1.4.1: target for configuration [library] is up to date.
  Up-to-date hunt-cache 0.10.1: target for configuration [library] is up to date.
  Up-to-date hunt-console 0.4.0: target for configuration [hunt-console] is up to date.
  Up-to-date hunt-sql 1.6.0: target for configuration [library] is up to date.
  Up-to-date hunt-database 2.3.6: target for configuration [default] is up to date.
  Up-to-date hunt-validation 0.5.0: target for configuration [library] is up to date.
  Up-to-date hunt-entity 2.8.1: target for configuration [library] is up to date.
  Up-to-date hunt-openssl 1.0.5: target for configuration [library] is up to date.
  Up-to-date hunt-jwt 0.2.0-beta.4: target for configuration [library] is up to date.
  Up-to-date hunt-security 0.6.0: target for configuration [library] is up to date.
  Up-to-date hunt-shiro 1.3.1: target for configuration [library] is up to date.
  Up-to-date poodinis 8.1.3: target for configuration [library] is up to date.
  Up-to-date hunt-framework 3.4.6: target for configuration [library] is up to date.
  Up-to-date webapp ~master: target for configuration [application] is up to date.
    Finished To force a rebuild of up-to-date targets, run again with --force
     Running webapp


 ___  ___     ___  ___     ________      _________
|\  \|\  \   |\  \|\  \   |\   ___  \   |\___   ___\     Hunt Framework 3.4.6
\ \  \\\  \  \ \  \\\  \  \ \  \\ \  \  \|___ \  \_|
 \ \   __  \  \ \  \\\  \  \ \  \\ \  \      \ \  \      Listening: 0.0.0.0:8080
  \ \  \ \  \  \ \  \\\  \  \ \  \\ \  \      \ \  \     TLS: Disabled
   \ \__\ \__\  \ \_______\  \ \__\\ \__\      \ \__\
    \|__|\|__|   \|_______|   \|__| \|__|       \|__|    https://www.huntframework.com


Try to browse http://0.0.0.0:8080
\end{minted}


\begin{minted}{text}
stud@sysrescue:~/projects/webapp/source/app/controller$ ps -FC webapp
UID          PID    PPID  C    SZ   RSS PSR STIME TTY          TIME CMD
stud       18551   18544  0 163131 37192  0 17:25 pts/1    00:00:00 /home/stud/projects/webapp/webapp
\end{minted}

\begin{minted}{bash}
stud@sysrescue:~/projects/webapp/source/app/controller$ curl http://localhost:8080/api/fact/9
{"result":"362880"}
stud@sysrescue:~/projects/webapp/source/app/controller$ curl http://localhost:8080/api/test
{"currtime":"2023-11-07T17:25:59.8807696"}
stud@sysrescue:~/projects/webapp/source/app/controller$ curl http://localhost:8080/api/echo/message-to-test
{"echo":"message-to-test"}
\end{minted}

\begin{minted}{html}
  <!--
$ stud@sysrescue:~/projects/webapp/source/app/controller$ curl http://localhost:8080/api/fact/9e
  -->
<!doctype html>
<html lang="en">
    <meta charset="utf-8">
    <title>404 Not Found</title>
    <meta name="viewport" content="width=device-width, initial-scale=1">
    <style>

        * {
            line-height: 3;
            margin: 0;
        }

        html {
            color: #888;
            display: table;
            font-family: sans-serif;
            height: 100%;
            text-align: center;
            width: 100%;
        }

        body {
            display: table-cell;
            vertical-align: middle;
            margin: 2em auto;
        }

        h1 {
            color: #555;
            font-size: 2em;
            font-weight: 400;
        }

        p {
            margin: 0 auto;
            width: 90%;
        }

    </style>
</head>
<body>
    <h1>404 Not Found</h1>
    <p>Sorry!! Unable to complete your request :(</p>

</body>
</html>
\end{minted}

\begin{minted}{d}
module app.controller.ApiController;

import hunt.framework;
import std.json : JSONValue;
import std.stdio;
import std.conv;

class ApiController : Controller
{
  mixin MakeController;

  // . . . . . . . . . . . . . .

  @Action
  JsonResponse fact(string n) {
    JSONValue js;
    auto res = fact(to!int(n));
    js["result"] = to!string(res);
    auto resp = new JsonResponse(js);
    return resp;
  }

  int fact(int n) {
    if (n==0) return 1;
    if (n==1) return 1;
    return n*fact(n-1);
  }
}
\end{minted}

\begin{minted}{llvm}
; [#uses = 1]
; Function Attrs: uwtable
define %hunt.framework.http.JsonResponse.JsonResponse* @_D3app10controller13ApiControllerQp4factMFAyaZC4hunt9framework4http12JsonResponseQo(%app.controller.ApiController.ApiController* nonnull %.this_arg, { i64, i8* } %n_arg) #0 {
  %this = alloca %app.controller.ApiController.ApiController*, align 8 ; [#uses = 3, size/byte = 8]
  %n = alloca { i64, i8* }, align 8               ; [#uses = 2, size/byte = 16]
  %js = alloca %std.json.JSONValue, align 8       ; [#uses = 3, size/byte = 24]
  %res = alloca i32, align 4                      ; [#uses = 2, size/byte = 4]
  %resp = alloca %hunt.framework.http.JsonResponse.JsonResponse*, align 8 ; [#uses = 3, size/byte = 8]
  store %app.controller.ApiController.ApiController* %.this_arg, %app.controller.ApiController.ApiController** %this, align 8
  store { i64, i8* } %n_arg, { i64, i8* }* %n, align 8
  %1 = bitcast %std.json.JSONValue* %js to i8*    ; [#uses = 1]
  call void @llvm.memcpy.p0i8.p0i8.i64(i8* align 1 %1, i8* align 1 bitcast (%std.json.JSONValue* @_D3std4json9JSONValue6__initZ to i8*), i64 24, i1 false)
  %2 = load %app.controller.ApiController.ApiController*, %app.controller.ApiController.ApiController** %this, align 8 ; [#uses = 1]
  %3 = getelementptr inbounds %app.controller.ApiController.ApiController, %app.controller.ApiController.ApiController* %2, i32 0, i32 0 ; [#uses = 1, type = [38 x i8*]**]
  %4 = load [38 x i8*]*, [38 x i8*]** %3, align 8 ; [#uses = 1]
  %"fact@vtbl" = getelementptr inbounds [38 x i8*], [38 x i8*]* %4, i32 0, i32 36 ; [#uses = 1, type = i8**]
  %5 = load i8*, i8** %"fact@vtbl", align 8       ; [#uses = 1]
  %fact = bitcast i8* %5 to i32 (%app.controller.ApiController.ApiController*, i32)* ; [#uses = 1]
  %6 = load %app.controller.ApiController.ApiController*, %app.controller.ApiController.ApiController** %this, align 8 ; [#uses = 1]
  %7 = load { i64, i8* }, { i64, i8* }* %n, align 8 ; [#uses = 1]
  %8 = call i32 @_D3std4conv__T2toTiZ__TQjTAyaZQqFNaNfQlZi({ i64, i8* } %7) #1 ; [#uses = 1]
  %9 = call i32 %fact(%app.controller.ApiController.ApiController* nonnull %6, i32 %8) ; [#uses = 1]
  store i32 %9, i32* %res, align 4
  %10 = load i32, i32* %res, align 4              ; [#uses = 1]
  %11 = call { i64, i8* } @_D3std4conv__T2toTAyaZ__TQlTiZQqFNaNbNfiZQx(i32 %10) #1 ; [#uses = 1]
  call void @_D3std4json9JSONValue__T13opIndexAssignTAyaZQuMFNaNfQmQoZv(%std.json.JSONValue* nonnull %js, { i64, i8* } %11, { i64, i8* } { i64 6, i8* getelementptr inbounds ([7 x i8], [7 x i8]* @.str.20, i32 0, i32 0) }) #1
  %.newclass_gc_alloc = call %object.Object* @_d_allocclass(%object.TypeInfo_Class* @_D4hunt9framework4http12JsonResponseQo7__ClassZ) #2 ; [#uses = 1]
  %.newclass_gc = bitcast %object.Object* %.newclass_gc_alloc to %hunt.framework.http.JsonResponse.JsonResponse* ; [#uses = 5]
  %vtbl = getelementptr inbounds %hunt.framework.http.JsonResponse.JsonResponse, %hunt.framework.http.JsonResponse.JsonResponse* %.newclass_gc, i32 0, i32 0 ; [#uses = 1, type = [19 x i8*]**]
  store [19 x i8*]* @_D4hunt9framework4http12JsonResponseQo6__vtblZ, [19 x i8*]** %vtbl, align 8
  %monitor = getelementptr inbounds %hunt.framework.http.JsonResponse.JsonResponse, %hunt.framework.http.JsonResponse.JsonResponse* %.newclass_gc, i32 0, i32 1 ; [#uses = 1, type = i8**]
  store i8* null, i8** %monitor, align 8
  %12 = getelementptr inbounds %hunt.framework.http.JsonResponse.JsonResponse, %hunt.framework.http.JsonResponse.JsonResponse* %.newclass_gc, i32 0, i32 2 ; [#uses = 1, type = %hunt.http.server.HttpServerResponse.HttpServerResponse**]
  %13 = getelementptr inbounds %hunt.framework.http.JsonResponse.JsonResponse, %hunt.framework.http.JsonResponse.JsonResponse* @_D4hunt9framework4http12JsonResponseQo6__initZ, i32 0, i32 2 ; [#uses = 1, type = %hunt.http.server.HttpServerResponse.HttpServerResponse**]
  %14 = bitcast %hunt.http.server.HttpServerResponse.HttpServerResponse** %12 to i8* ; [#uses = 1]
  %15 = bitcast %hunt.http.server.HttpServerResponse.HttpServerResponse** %13 to i8* ; [#uses = 1]
  call void @llvm.memcpy.p0i8.p0i8.i64(i8* align 1 %14, i8* align 1 %15, i64 9, i1 false)
  %16 = call %hunt.framework.http.JsonResponse.JsonResponse* @_D4hunt9framework4http12JsonResponseQo__T6__ctorTS3std4json9JSONValueZQBdMFQBaZCQDaQCyQCrQCpQCs(%hunt.framework.http.JsonResponse.JsonResponse* nonnull returned %.newclass_gc, %std.json.JSONValue* byval(%std.json.JSONValue) align 8 %js) #1 ; [#uses = 0]
  store %hunt.framework.http.JsonResponse.JsonResponse* %.newclass_gc, %hunt.framework.http.JsonResponse.JsonResponse** %resp, align 8
  %17 = load %hunt.framework.http.JsonResponse.JsonResponse*, %hunt.framework.http.JsonResponse.JsonResponse** %resp, align 8 ; [#uses = 0]
  %18 = load %hunt.framework.http.JsonResponse.JsonResponse*, %hunt.framework.http.JsonResponse.JsonResponse** %resp, align 8 ; [#uses = 1]
  ret %hunt.framework.http.JsonResponse.JsonResponse* %18
}
\end{minted}


\chapter{Реализация ....}

\section{Скрипт порождения структуры базы данных}
\label{sec:struct-bd}

Начало абзаца


\begin{figure}[htbp]
\begin{center}
\begin{minted}[fontsize=\normalsize]{sql}
CREATE TABLE Persons (
    PersonID int,
    LastName varchar(255),
    FirstName varchar(255),
    Address varchar(255),
    City varchar(255)
);
\end{minted}
\end{center}
\caption{Скрипт SQL для создания таблицы Persons}\label{fig:create-table}
\end{figure}

\begin{lstlisting}[
  caption={Скрипт SQL для создания таблицы Persons},  % Подпись
  captionpos=b,               % Подпись снизу
  basicstyle=\small\ttfamily, % Размер и тип шрифта для представления листинга
  language=SQL
  ]
CREATE TABLE Persons (
    PersonID int,
    LastName varchar(255),
    FirstName varchar(255),
    Address varchar(255),
    City varchar(255)
);
\end{lstlisting}


\chapter{Тестирование программной системы в нагруженной среде}

Разработан набор тестов структуры БД (см.~\ref{sec:struct-bd} на стр.~\pageref{sec:struct-bd}).   Загрузка скрипта базы данных осуществляется при помощи команды \texttt{mysql}.  Дамп существующей БД выполняется командой \verb|mysql-dump|.

Это тест шрифта \textsf{без засечек, А}А.
Это тест шрифта \texttt{моноширинный, А}А.
Это тест шрифта \textit{наклонный, А}А.

\begin{figure}[hbtp]
  \centering
  \includegraphics[width=0.3\linewidth]{fig1.pdf}
  \caption{Пример работы программы}
  \label{fig:prog-ex}
\end{figure}


\begin{equation}
  \label{eq:f1}
  E=mc^2,
\end{equation}

\noindent{}где $m$ -- это масса, $c$ -- скорость света в вакууме, \ldots

Новый абзац. Рассмотрим формулу (\ref{eq:f1}).


\chapter*{ЗАКЛЮЧЕНИЕ}

В таблице~\ref{tbl:ex} представлены результаты сравнения производительности\ldots

\begin{table}[hb]
  \caption{Пример таблицы}\label{tbl:ex}
  \centering
\begin{tabularx}{1\textwidth} {
  | >{\raggedright\arraybackslash}X
  | >{\centering\arraybackslash}X
  | >{\raggedleft\arraybackslash}X | }
 \hline
 item 11 & item 12 & item 13 \\
 \hline
 item 21  & item 22  & item 23  \\
\hline
\end{tabularx}
\end{table}

\begin{thebibliography}{99}
\bibitem{bratko90} И.~Братко. Язык программирования Пролог для искусственного интеллекта. М\;:~Наука. 1990. 310~c.
\bibitem{b2} DeRidder J.L. The immediate prospects for the application of ontologies in digital libraries // Knowledge Organization -- 2007. -- Vol. 34, No. 4. P.~227-246.
\bibitem{b2}  U.S. National Library of Medicine. Fact sheet: UMLS Metathesaurus\;:\;[текст]~/ National Institutes of Health, 2006-2013. -- URL:
\url{http://www.nlm.nih.gov/pubs/factsheets/umlsmeta.html} (дата обращения: 2014-12-09).
\bibitem{b2}  U.S. National Library of Medicine. Fact sheet: Unfied Medical Language System\;:\;[текст]~/ National Institutes of Health, 2006--2013. -- URL:\url{http://www.nlm.nih.gov/pubs/factsheets/umls.html} (дата обращения: 2009-12-09).
\bibitem{b3}  Антопольский А.Б., Белоозеров В.Н. Процедура формирования макротезауруса политематических информационных систем // Классификация и кодирование. -- 1976. -- N 1 (57). -- С.~25--29.
\bibitem{b4}  Белоозеров В.Н., Федосимов В.И. Место макротезауруса в лингвистическом обеспечении сети органов научнотехнической информации // Проблемы информационных систем. -- 1986. -- N 1. -- С.~6--10.
\bibitem{b5}  Использование и ведение макротезауруса ГАСНТИ:Методические рекомендации / ГКНТСССР -- М., 1983. -- 12~с.
\bibitem{b6}  Nuovo soggettario: guidaalsistemaitaliano di indicizzazione per soggetto, prototipo del thesaurus\;:\;[рецензия]~// Knowledge Organization. -- 2007. -- Vol. 34, N 1. -- P.~58--60.
\bibitem{b7}  ГОСТ 7.25-2001 СИБИД. Тезаурус информационно-поисковый одноязычный. Правила разработки, структура, состав и форма представления. -- М., 2002. -- 16~с.
\bibitem{b8}  Nanoscale Science and Technology Supplement: Collection ofapplicable terms fromPACS 2008\;:\;[текст]~// PACS 2010 Regular Eddition /~AIP Publishing. -- URL: \url{http://www.aip.org/publishing/pacs/nano-supplement} (дата обращения: 2014-12-09).
\bibitem{b9}  Смирнова О.В. Методикасоставления индексов УДК // Научно-техническая информация. Сер.~1. -- 2008. -- N 8. -- С.~7--8.
\bibitem{b10}  Индексирование фундаментальных научных направлений кодами информационных классификаций УДК /~О.А.~Антошкова, Т.С.~Астахова, В.Н.~Белоозеров и др.; под ред.акад. Ю.М.~Арского. -- М., 2010. -- 322~с.
\bibitem{b11}  Рубрикатор как инструмент информационной навигации /~Р.С.~Гиляревский, А.В.~Шапкин, В.Н.~Белоозеров. - СПб.\;: Профессия, 2008. -- 352~с.
\bibitem{b12}  Рубрикатор научно-технической информации по нанотехнологиям и наноматериалам / РНЦ "Курчатовский институт",
ФГУ ГНИИ ИТТ "Информика", Национальный электронно-информационный консорциум (НЭИКОН), Всероссийский институт научной и технической информации (ВИНИТИРАН). -- М., 2009. -- 75~с.
\bibitem{b13}  Рубрикатор по нанонауке и нанотехнологиям\;:\;[сайт]~-- URL:\url{http/www.rubric.neicon.ru} (дата обращения: 12-02-2022)
\end{thebibliography}

\appendices

\chapter{Исходный код программ}
\chapter{Документация разработчика}

\end{document}




%%% Local Variables:
%%% mode: latex
%%% TeX-master: t
%%% End:
